\chapter{PENDAHULUAN}
\hyphenpenalty=10000
\exhyphenpenalty=10000
\sloppy

\section{Deskripsi Aplikasi}
Transformasi digital telah menjadi kebutuhan strategis bagi institusi pendidikan tinggi dalam meningkatkan efisiensi operasional dan kualitas layanan akademik \parencite{Ramdany2024}. Konsep \textit{smart campus} yang memanfaatkan teknologi informasi seperti \textit{cloud computing}, \textit{big data}, dan infrastruktur jaringan terintegrasi, menawarkan peluang untuk mengoptimalkan berbagai aspek pengelolaan kampus \parencite{Nachandiya2018}. Sarana dan prasarana merupakan komponen pendukung utama kegiatan akademik dan non-akademik, sehingga pengelolaannya perlu dilakukan secara terencana dan terdokumentasi. Namun, pada banyak institusi pendidikan, pengelolaan dan peminjaman fasilitas masih dilakukan secara manual atau menggunakan media pencatatan terpisah seperti \textit{Microsoft Excel}, yang berpotensi menimbulkan kesulitan pencarian data, keterlambatan laporan, serta rendahnya efisiensi kerja pengelola \parencite{PutriNabella2025}.

Perkembangan teknologi informasi mendorong penerapan sistem informasi berbasis web sebagai solusi atas permasalahan tersebut. Sistem berbasis web memungkinkan pengelolaan data secara terpusat, mudah diakses, serta mendukung proses pengajuan, pencatatan, dan pelaporan peminjaman fasilitas secara lebih efektif dan terstruktur \parencite{Purwati2022}. Penyajian informasi peminjaman secara \textit{real-time} dapat meminimalkan kesalahan informasi dan bentrok jadwal penggunaan fasilitas kampus \parencite{Herlambang2020}. Integrasi data dan otomatisasi proses dalam sistem informasi juga terbukti memberikan manfaat signifikan terhadap efisiensi operasional dan peningkatan kualitas layanan \parencite{Ramdany2024}.

Berdasarkan kondisi tersebut, dikembangkan Sistem Informasi Peminjaman Sarana dan Prasarana Kampus Berbasis Web menggunakan CRUD RDBMS sebagai solusi terintegrasi dalam pengelolaan peminjaman fasilitas. Sistem ini dirancang untuk mendukung proses pengajuan, persetujuan, serta pendokumentasian peminjaman secara efektif, transparan, dan akuntabel.

\newpage

\section{Rumusan Masalah}
Berdasarkan latar belakang yang telah diuraikan sebelumnya, maka permasalahan dalam Proyek II ini dapat dirumuskan dalam bentuk pertanyaan sebagai berikut:
\begin{enumerate}
    \item Bagaimana merancang sistem informasi berbasis web yang mampu mengelola proses peminjaman sarana dan prasarana kampus secara efektif, terintegrasi, dan terkomputerisasi?
    \item Bagaimana menyediakan mekanisme pengajuan dan pemantauan status peminjaman sarana dan prasarana kampus yang mudah diakses oleh mahasiswa secara real-time?
    \item Bagaimana sistem informasi dapat meminimalkan terjadinya bentrok jadwal peminjaman serta mendukung pengelolaan data peminjaman yang terdokumentasi dengan baik untuk keperluan pelaporan dan evaluasi?
\end{enumerate}

\section{Tujuan}
Berdasarkan rumusan masalah yang telah diuraikan sebelumnya, maka tujuan dari pelaksanaan Proyek II ini adalah sebagai berikut:
\begin{enumerate}
    \item Untuk merancang dan membangun sistem informasi peminjaman sarana dan prasarana kampus berbasis web yang terintegrasi guna mendukung proses pengelolaan fasilitas kampus secara efektif dan terstruktur.
    \item Untuk menyediakan mekanisme pengajuan serta pemantauan status peminjaman sarana dan prasarana kampus yang mudah diakses oleh mahasiswa secara real-time.
    \item Untuk meminimalkan terjadinya bentrok jadwal peminjaman melalui fitur validasi jadwal serta menyediakan pengelolaan data peminjaman yang terdokumentasi dengan baik guna mendukung keperluan pelaporan dan evaluasi.
\end{enumerate}

\newpage

\section{Lingkup Dokumentasi}
Ruang lingkup dokumentasi dalam laporan ini meliputi:
\begin{enumerate}
    \item Perancangan dan implementasi sistem peminjaman sarana dan prasarana kampus yang mencakup proses pengajuan, persetujuan, dan pencatatan kehadiran peminjam.
    \item Analisis kebutuhan fungsional dan non-fungsional sistem yang mendukung pengelolaan peminjaman secara terintegrasi.
    \item Pengelolaan data inventaris sarana dan prasarana yang dibatasi pada operasi tambah, ubah (edit), dan hapus data, tanpa mencakup proses pengadaan atau manajemen stok lanjutan.
\end{enumerate}
